\documentclass[12pt]{article}
\usepackage[margin=1in]{geometry}
\usepackage{setspace}
\onehalfspacing

% Start of preamble
%==========================================================================================%
% Required to support mathematical unicode
\usepackage[warnunknown, fasterrors, mathletters]{ucs}
\usepackage[utf8x]{inputenc}

\usepackage[dvipsnames,table,xcdraw]{xcolor}
\usepackage{hyperref} 
\hypersetup{
colorlinks=true,
linkcolor=blue,
filecolor=magenta,
urlcolor=cyan,
pdfpagemode=FullScreen
}

% Standard mathematical typesetting packages
\usepackage{amsmath,amssymb,amscd,amsthm,amsxtra, pxfonts}
\usepackage{mathtools,mathrsfs,dsfont,xparse}

% Symbol and utility packages
\usepackage{cancel, textcomp}
\usepackage[mathscr]{euscript}
\usepackage[nointegrals]{wasysym}
\usepackage{apacite}

% Extras
\usepackage{physics}  
\usepackage{tikz-cd} 
\usepackage{microtype}
\usepackage{enumitem}
\usepackage{titling}
\usepackage{graphicx}

% Fancy theorems due to @intuitively on discord
\usepackage{mdframed}
\newmdtheoremenv[
backgroundcolor=NavyBlue!30,
linewidth=2pt,
linecolor=NavyBlue,
topline=false,
bottomline=false,
rightline=false,
innertopmargin=10pt,
innerbottommargin=10pt,
innerrightmargin=10pt,
innerleftmargin=10pt,
skipabove=\baselineskip,
skipbelow=\baselineskip
]{mytheorem}{Theorem}

\newenvironment{theorem}{\begin{mytheorem}}{\end{mytheorem}}

\newtheorem{corollary}{Corollary}
\newtheorem{lemma}{Lemma}

\newtheoremstyle{definitionstyle}
{\topsep}%
{\topsep}%
{}%
{}%
{\bfseries}%
{.}%
{.5em}%
{}%
\theoremstyle{definitionstyle}
\newmdtheoremenv[
backgroundcolor=Violet!30,
linewidth=2pt,
linecolor=Violet,
topline=false,
bottomline=false,
rightline=false,
innertopmargin=10pt,
innerbottommargin=10pt,
innerrightmargin=10pt,
innerleftmargin=10pt,
skipabove=\baselineskip,
skipbelow=\baselineskip,
]{mydef}{Definition}
\newenvironment{definition}{\begin{mydef}}{\end{mydef}}

\newtheorem*{remark}{Remark}

\newtheorem*{example}{Example}

% Common shortcuts
\def\mbb#1{\mathbb{#1}}
\def\mfk#1{\mathfrak{#1}}

\def\bN{\mbb{N}}
\def \C{\mbb{C}}
\def \R{\mbb{R}}
\def\bQ{\mbb{Q}}
\def\bZ{\mbb{Z}}
\def \cph{\varphi}
\renewcommand{\th}{\theta}
\def \ve{\varepsilon}
\newcommand{\mg}[1]{\| #1 \|}

% Often helpful macros
\newcommand{\floor}[1]{\left\lfloor#1\right\rfloor}
\newcommand{\ceil}[1]{\left\lceil#1\right\rceil}
\renewcommand{\qed}{\hfill\qedsymbol}
\renewcommand{\P}{\mathbb P\qty}
\newcommand{\E}{\mathbb{E}\qty}
\newcommand{\Cov}{\mathrm{Cov}\qty}
\newcommand{\Var}{\mathrm{Var}\qty}

% Sets
\usepackage{braket}

\graphicspath{{/}}
\usepackage{float}

\newcommand{\SET}[1]{\Set{\mskip-\medmuskip #1 \mskip-\medmuskip}}

% End of preamble
%==========================================================================================%

% Start of commands specific to this file
%==========================================================================================%

%==========================================================================================%
% End of commands specific to this file

\title{CSE Template}
\date{\today}
\author{Rohan Mukherjee}

\begin{document}
    \maketitle
    \subsection*{Problem 1.}
    We use a truly mindblowing trick to solve this problem. First replace every edge of $G$ with two directed edges, one going forward and the other going backwards. Then for each $v \in G$, replace $v$ with two vertices $v_{in}$ and $v_{out}$, with one edge $v_{in} \to v_{out}$, and for each neighbor edge $w \to v$ add $w \to v_in$ and for each $v \to u$, add $v_{out} \to u$. After running this on the entire graph, let $\hat S$ be the collection of $s_{in}$'s attained in this way from the original vertices in $S$, and similarly $\hat T$ to be  the collection of $t_{in}$'s attained in this way from the original vertices in $T$. Remove all the $t_{out}$s and $s_{in}$s, and add the vertices $\hat S$ and $\hat T$ with edge $\hat S \to s_{out}$ for each $s_{out}$ in the graph, and $t_{in} \to \hat T$. Return the number of edge-disjoint paths from $\hat S$ to $\hat T$ in the new graph.

    We now prove the correctness of the above algorithm. Suppose we start with $\mathscr{P}$ being some set of edge disjoint paths in the original graph. Let $P = s = v^{(1)}, \ldots, v^{(n)}= t$ be one of those paths. We then can put $P$ on the transformed graph $G'$ by $v^{(1)}_{in} \to v^{(1)}_{out}, \ldots, v^{(n)}_{in} \to v^{(n)}_{out}$. This is a path in the new graph by construction, and the set of such paths are edge-disjoint because the original ones were vertex disjoint by what we did in class. 

    Similarly, given a set of edge-disjoint paths in the new graph, we can put them backe on the original graph by noticing the following: if $v_{in}$ is in the path $P$ for some $v \in G$, then we must have $v_{out}$ be the next vertex in the path because that is the only place $v_{in}$ goes to. Recalling then that we start at some $s_{in}$ for some $s \in S$ and end at some $t_{out}$ for some $t \in T$, we then notice that every path is of the following type:
    \begin{align*}
        s_{in} \to s_{out} \to v_{in} \to v_{out} \to \ldots \to t_{in} \to t_{out}
    \end{align*}
    From here we can put this path on the original graph by decreeing that $s \to v \to \ldots \to t$ in the original graph. The set of such paths in the original graph is vertex disjoint because otherwise we would be using the edge $v_{in} \to v_{out}$ multiple times, which completes the proof.

    \subsection*{Problem 2.}
    Suppose instead that the max number of edge-disjoint paths from $s$ to $u$ is $\ell < k$. By Mengers theorem, this says that the min cut between $s$ and $u$ is just $\ell$. Thus we can remove $\ell$ edges from the digraph $G$ to separate $s$ and $u$, i.e. put them into two different connected components. Since there are at least $k$ edge-disjoint paths from $s$ to $t$ and $t$ to $u$, we need to remove at least $k$ edges to separate $s$ from $t$ or $t$ from $u$. In particular, since we removed $<k$ edges from the graph, there is still at least one path from $s \to t$ and $t \to u$. Let such two paths be $P_1$ and $P_2$. By looking at the last time $P_2$ intersects $P_1$, and removing all edges before that, we get a path from $s \to u$, a contradiction, because by hypothesis $s$ and $u$ are separated.
\end{document}